\subsection{模型建立}
在这一问里，我们把长期动力学部分压缩为 Hastings-Powell 三级食物链这一最小闭环，在同一参数轴上比较“资源--消费者--捕食者”反馈下的轨道类型\cite{hp1991}：
\begin{equation}
\frac{\mathrm{d}X_1}{\mathrm{d}t}
=rX_1\left(1-\frac{X_1}{K}\right)-\frac{a_1X_1X_2}{b_1+X_1},
\end{equation}
\begin{equation}
\frac{\mathrm{d}X_2}{\mathrm{d}t}
=\frac{a_1X_1X_2}{b_1+X_1}-\frac{a_2X_2X_3}{b_2+X_2}-d_1X_2,
\end{equation}
\begin{equation}
\frac{\mathrm{d}X_3}{\mathrm{d}t}
=\frac{a_2X_2X_3}{b_2+X_2}-d_2X_3.
\end{equation}
其中 $K$ 代表基础食物链有效承载力，可理解为营养输入、水质改善或污染冲击后的综合结果。物种映射上，$X_1$ 对应底层资源层（问题一中 $R_j$ 的聚合），$X_2$ 对应中层消费者（四大家鱼功能群），$X_3$ 对应顶层捕食者（江豚等）。我们在这里仅保留营养级对应，用于说明三级食物链的生态位含义，不直接继承问题一的参数数值。

模型参数取值如表 \ref{tab:q3_hp_params} 所示，其中 $a_1,b_1,d_1$ 和 $a_2,b_2,d_2$ 分别控制中层和顶层的摄食与损失，$r$ 固定为 1，$K$ 作为扫描控制参数。
\begin{table}[htbp]
\centering
\caption{问题三的 Hastings-Powell 模型参数取值}
\label{tab:q3_hp_params}
\begin{tabularx}{0.85\textwidth}{>{\raggedright\arraybackslash}p{1.5cm} >{\centering\arraybackslash}p{1.5cm} X}
\toprule
参数 & 取值 & 说明 \\
\midrule
$r$ & 1.0 & 底层资源内禀增长率，固定为归一化基准 \\
$a_1$ & 5.0 & 中层对底层的摄食率，取自 Hastings \& Powell (1991) 默认值\cite{hp1991} \\
$b_1$ & 1.0 & 中层摄食半饱和常数 \\
$a_2$ & 0.1 & 顶层对中层的摄食率 \\
$b_2$ & 2.0 & 顶层摄食半饱和常数 \\
$d_1$ & 0.4 & 中层自然死亡率 \\
$d_2$ & 0.01 & 顶层自然死亡率 \\
$K$ & 扫描 & 承载力控制参数，扫描范围 $[1.4,4.8]$ \\
\bottomrule
\end{tabularx}
\end{table}

